% !TeX root = booklet-03.tex
% Reusable three-terminal symbols following the source figures.
\tikzset{
 pics/study-npass/.style={code={
  \draw (0,0)--(0.25,0)--(0.25,0.2)--(0.75,0.2)--(0.75,0)--(1,0);
  \draw (0.15,0.36)--(0.85,0.36);
  \draw (0.5,0.36)--(0.5,0.85);
  \coordinate (-L) at (0,0);
  \coordinate (-R) at (1,0);
  \coordinate (-G) at (0.5,0.85);
 }},
 pics/study-ppass/.style={code={
  \draw (0,0)--(0.25,0)--(0.25,-0.2)--(0.75,-0.2)--(0.75,0)--(1,0);
  \draw (0.15,-0.36)--(0.85,-0.36);
  \draw (0.5,-0.44) circle[radius=0.065];
  \draw (0.5,-0.505)--(0.5,-0.85);
  \coordinate (-L) at (0,0);
  \coordinate (-R) at (1,0);
  \coordinate (-G) at (0.5,-0.85);
 }}
}
\clearpage
\subsection{مشخصه \lr{DC} یک \lr{Pass Transistor}}
\label{sec:pass-transistor}
در \lr{Pass Transistor}، سیگنال از مسیر کانال بین دو پایانه
عبور می‌کند و \lr{Gate} روشن یا خاموش‌بودن مسیر را کنترل
می‌کند. این ساختار برای انتقال یک ولتاژ به کار می‌رود.
تحلیل این بخش بر صفحات ۹۲ و آغاز ۹۳ کتاب، شکل
\secnum{۲.۳۱} و اسلایدهای ۱۲ و ۱۳ تکیه دارد
\cite{weste,dcslides}. حل تمرین \secnum{۲.۲۰} با شکل اصلی
\secnum{۲.۳۵} در صفحه چاپی ۹۷ تطبیق داده شده است.

\subsubsection{پایانه‌ها، فرض بار و معنای حد انتقال}
\label{sec:pass-conventions}
در مدل دوطرفه کانال که اینجا استفاده می‌کنیم، برای \lr{nMOS}
پایانه دارای ولتاژ کمتر نقش \lr{Source} و برای \lr{pMOS}
پایانه دارای ولتاژ بیشتر نقش \lr{Source} را در روابط جریان
دارد. بنابراین، محل قرارگرفتن پایانه در سمت چپ یا راست شکل
به‌تنهایی نقش آن را تعیین نمی‌کند. \lr{Gate} پایانه کنترل است
و باید اتصال آن را جدا از مسیر عبور سیگنال دنبال کرد.
بدنه \lr{nMOS} معمول را زمین و بدنه \lr{pMOS} را تغذیه
می‌گیریم، مگر آنکه فرض دیگری صریحاً نوشته شود.

ابتدا آستانه‌های ثابت \(V_{tn}>0\) و
\(V_{TP}=|V_{tp}|>0\)، بار بدون جریان \lr{DC} و مدل بدون
نشتی را به کار می‌بریم. ولتاژها در بازه زمین تا تغذیه‌اند.
خروجی می‌تواند ظرفیت داشته باشد. در چنین مدلی، جریان
ترانزیستور، بار ظرفیت را تغییر می‌دهد تا اختلاف دو پایانه صفر
شود یا افزاره به مرز خاموشی برسد.

اگر خروجی با رسیدن به مرز آستانه خاموش شود، مدل ایده‌آل
امکان نگهداری بار را دارد. پس حد شارژ از پایین یا حد دشارژ
از بالا باید با مسیر حرکت گره بیان شود. برای یک گره شناور
با شرایط اولیه دلخواه، معادله جریان صفر می‌تواند چندین
ولتاژ خروجی را بپذیرد. کتاب نیز در پانوشت صفحه ۹۲، امکان
افزایش بسیار آهسته ولتاژ خروجی \lr{nMOS} بر اثر
\lr{Subthreshold Leakage} را یادآوری می‌کند \cite{weste}.

\subsubsection{عبور سطح بالا و پایین با آستانه ثابت}
\label{sec:pass-threshold-drop}
فرض کنید ورودی و \lr{Gate} یک \lr{nMOS} به
\(V_{DD}\) وصل باشند و خروجی از صفر شارژ شود.
در طول شارژ، خروجی پایانه کم‌ولتاژتر و نقش \lr{Source} است.
با \(y=V_{\mathrm{out}}\)، اضافه‌ولتاژ برابر است با
\begin{equation}
 V_{GS}-V_{tn}=V_{DD}-y-V_{tn}.
 \label{eq:pass-n-overdrive}
\end{equation}
با افزایش \(y\)، این اضافه‌ولتاژ کم می‌شود. اگر
\(V_{DD}>V_{tn}\)، حد هدایت قوی با صفرشدن آن به دست می‌آید
\begin{equation}
 \boxed{y_{n,\mathrm{lim}}=V_{DD}-V_{tn}}.
 \label{eq:pass-n-high-limit}
\end{equation}
این کاهش سطح بالا نسبت به ورودی، \lr{Threshold Drop} نامیده
می‌شود \cite[شکل \secnum{۲.۳۱}\lr{(a)}]{weste}.
برای \lr{Gate} ثابت \(V_G>V_{tn}\)، ورودی ثابت
\(v\ge0\) و شارژ خروجی از صفر، نتیجه مدل ثابت‌آستانه
به‌صورت زیر تعمیم پیدا می‌کند
\begin{equation}
 y_{n,\mathrm{lim}}=\min(v,V_G-V_{tn}).
 \label{eq:pass-n-min-rule}
\end{equation}
اگر ورودی پایین‌تر از حد گیت باشد، خروجی می‌تواند به خود ورودی
برسد. اگر ورودی بالاتر باشد، افزاره هنگام شارژ به مرز آستانه
می‌رسد. این رابطه به فرض شارژ از پایین و شرایط روشن‌شدن وابسته است.

وقتی همین \lr{nMOS} با \lr{Gate} بالا خروجی را به زمین
متصل می‌کند، پایانه زمین نقش \lr{Source} دارد و
\(V_{GS}=V_G\) هنگام دشارژ کاهش نمی‌یابد.
برای \(V_G>V_{tn}\)، مسیر تا رسیدن خروجی به صفر هدایت دارد.
در حد نهایی، اختلاف درین و سورس صفر و جریان نیز صفر می‌شود.
ازاین‌رو، \lr{nMOS} سطح پایین را بدون افت آستانه منتقل می‌کند.

برای \lr{pMOS} با \lr{Gate} صفر و ورودی زمین، خروجی در
دشارژ از بالا نقش \lr{Source} دارد. شرط روشن‌بودن چنین است
\begin{equation}
 V_{SG}>V_{TP}\quad\Longrightarrow\quad y>V_{TP}.
 \label{eq:pass-p-overdrive}
\end{equation}
پس حد دشارژ با آستانه ثابت برابر است با
\begin{equation}
 \boxed{y_{p,\mathrm{lim}}=V_{TP}}.
 \label{eq:pass-p-low-limit}
\end{equation}
برای \lr{Gate} ثابت \(V_G\)، ورودی \(v\) و دشارژ از
خروجی بالا، در صورت وجود هدایت اولیه، قاعده متناظر می‌شود
\begin{equation}
 y_{p,\mathrm{lim}}=\max(v,V_G+V_{TP}).
 \label{eq:pass-p-max-rule}
\end{equation}
در مقابل، برای ورودی \(V_{DD}\) و \lr{Gate} صفر، پایانه
ورودی نقش \lr{Source} دارد. \(V_{SG}=V_{DD}\) ثابت می‌ماند
و با \(V_{DD}>V_{TP}\)، خروجی می‌تواند تا خود تغذیه شارژ شود.
جدول \ref{tab:pass-basic-levels} نتیجه این چهار حالت را
جمع می‌کند \cite[صفحات ۹۲ و ۹۳]{weste}.
\begin{table}[!htbp]
\centering
\small
\caption{انتقال دو ریل در مدل بدون اثر بدنه و بدون بار \lr{DC}.
خروجی از سمت دارای مسیر هدایت به حد نهایی نزدیک می‌شود.}
\label{tab:pass-basic-levels}
\begin{tabular}{cccc}
\toprule
افزاره & \lr{Gate} & ورودی & خروجی یا حد انتقال\\
\midrule
\lr{nMOS} & \(V_{DD}\) & صفر & صفر\\
\lr{nMOS} & \(V_{DD}\) & \(V_{DD}\) & \(V_{DD}-V_{tn}\)\\
\lr{pMOS} & صفر & صفر & \(V_{TP}\)\\
\lr{pMOS} & صفر & \(V_{DD}\) & \(V_{DD}\)\\
\bottomrule
\end{tabular}
\end{table}

\paragraph*{مثال آموزشی قاعده انتقال}
با فرض \(V_G=V_{DD}=\qty{1.2}{\volt}\) و
\(V_{tn}=\qty{0.3}{\volt}\)، حد شارژ \lr{nMOS} از پایین
\(\qty{0.9}{\volt}\) است. ورودی
\(\qty{0.6}{\volt}\) تا همان \(\qty{0.6}{\volt}\) منتقل
می‌شود. ورودی \(\qty{1.2}{\volt}\)، خروجی را در مدل به حد
\(\qty{0.9}{\volt}\) می‌رساند.
برای \lr{pMOS} با \lr{Gate} صفر و
\(V_{TP}=\qty{0.3}{\volt}\)، دشارژ به سمت ورودی صفر در
حد \(\qty{0.3}{\volt}\) متوقف می‌شود. اگر ورودی
\(\qty{0.6}{\volt}\) باشد، همان سطح بالاتر از حد گیت
قابل انتقال است.

\subsubsection{اثر بدنه و حل معادله حد خروجی}
\label{sec:pass-body-effect}
در \lr{nMOS} عبوردهنده سطح بالا با بدنه زمین،
\(V_{SB}=y\) است. افزایش خروجی، آستانه را نیز بیشتر می‌کند.
مدل بخش \secnum{۲.۴.۳.۱} کتاب را دوباره می‌نویسیم
\begin{equation}
 V_{tn}(y)=V_{tn0}+\gamma_n
       \left(\sqrt{\phi_n+y}-\sqrt{\phi_n}\right).
 \label{eq:pass-body-threshold}
\end{equation}
در اینجا \(V_{tn0}>0\) آستانه در بایاس بدنه صفر،
\(\phi_n>0\) پتانسیل سطحی با واحد ولت و
\(\gamma_n>0\) ضریب اثر بدنه با واحد
\(\unit{\volt}^{1/2}\) است.
حاصل‌ضرب \(\gamma_n\) در تفاضل ریشه‌ها واحد ولت دارد.
برای ورودی بالا و \lr{Gate} در \(V_G>V_{tn0}\)، حد شارژ
از رابطه زیر تعیین می‌شود
\begin{equation}
 y=V_G-V_{tn0}-\gamma_n
       \left(\sqrt{\phi_n+y}-\sqrt{\phi_n}\right).
 \label{eq:pass-body-implicit}
\end{equation}
آستانه به خروجی مجهول وابسته است، پس باید همین معادله حل شود.

برای حذف ریشه، \(z=\sqrt{\phi_n+y}\) تعریف می‌کنیم.
در نتیجه، \(y=z^2-\phi_n\) و برای خروجی نامنفی،
\(z\ge\sqrt{\phi_n}\) است. جای‌گذاری و جمع جمله‌ها می‌دهد
\begin{align}
 z^2-\phi_n&=V_G-V_{tn0}-\gamma_n z+\gamma_n\sqrt{\phi_n},\notag\\
 z^2+\gamma_nz
  -\left(\phi_n+V_G-V_{tn0}+\gamma_n\sqrt{\phi_n}\right)&=0.
 \label{eq:pass-body-quadratic}
\end{align}
فرمول درجه دوم، ریشه فیزیکی را به صورت زیر می‌دهد
\begin{equation}
 z=\frac{-\gamma_n+
 \sqrt{\gamma_n^2+4(\phi_n+V_G-V_{tn0}+\gamma_n\sqrt{\phi_n})}}{2},
 \qquad y=z^2-\phi_n.
 \label{eq:pass-body-root}
\end{equation}
ریشه با علامت منفی جلوی رادیکال، \(z<0\) می‌دهد و با
تعریف ریشه نامنفی سازگار نیست. همه جمله‌های داخل رادیکال
واحد ولت دارند و \(z\) واحد ریشه ولت دارد.

برای کنترل یکتایی، تابعی با واحد ولت تعریف می‌کنیم
\begin{equation}
 f(y)=y+V_{tn0}+\gamma_n
 (\sqrt{\phi_n+y}-\sqrt{\phi_n})-V_G.
 \label{eq:pass-body-residual}
\end{equation}
مشتق در دامنه خروجی نامنفی مثبت است.
\[
 f'(y)=1+\frac{\gamma_n}{2\sqrt{\phi_n+y}}>0,\qquad y\ge0.
\]
همچنین \(f(0)=V_{tn0}-V_G<0\) و
\(f(V_G-V_{tn0})>0\). پس یک ریشه در
\(0<y<V_G-V_{tn0}\) وجود دارد. اثر بدنه حد بالا را کاهش
داده است. در حد \(\gamma_n\to0\)، رابطه آستانه ثابت
\(y=V_G-V_{tn0}\) بازیابی می‌شود.

\paragraph*{مثال آموزشی و تفسیر عدد مرجع}
برای یک محاسبه آموزشی، فرض کنید
\(V_G=\qty{1}{\volt}\)، \(V_{tn0}=\qty{0.3}{\volt}\)،
\(\phi_n=\qty{0.93}{\volt}\) و
\(\gamma_n=0.16\,\unit{\volt}^{1/2}\).
رابطه \eqref{eq:pass-body-root} می‌دهد
\begin{equation}
 y\approx\qty{0.6530}{\volt},\qquad
 V_{tn}(y)\approx\qty{0.3470}{\volt}.
 \label{eq:pass-body-example}
\end{equation}
کنترل \(y+V_{tn}(y)\approx\qty{1}{\volt}\) و
\(0<y<\qty{0.7}{\volt}\) برقرار است.
کتاب برای نمونه فناوری خود در تغذیه \(\qty{1}{\volt}\)،
خروجی حدود \(\qty{0.65}{\volt}\) را بیان می‌کند
\cite[صفحه ۹۲]{weste}. پارامترهای گرد این مثال آموزشی
صریحاً فرض شده‌اند و عدد دقیق یک فرایند را مشخص نمی‌کنند.

برای \lr{pMOS} با بدنه \(V_{DD}\) و دشارژ خروجی به زمین،
\(V_{BS}=V_{DD}-y\) است. بزرگی آستانه و حد دشارژ در
\lr{Gate} صفر چنین به هم مربوط می‌شوند
\begin{equation}
 y=V_{TP0}+\gamma_p
   \left(\sqrt{\phi_p+V_{DD}-y}-\sqrt{\phi_p}\right).
 \label{eq:pass-p-body-implicit}
\end{equation}
در این رابطه، \(V_{TP0},\phi_p,\gamma_p\) همتای پارامترهای
\lr{nMOS} هستند. کاهش خروجی، بایاس معکوس سورس و بدنه را
بیشتر و بزرگی آستانه را افزایش می‌دهد.
با پارامترهای متقارن مثال قبل و \(V_{DD}=\qty{1}{\volt}\)،
تقارن دو رابطه می‌دهد
\[
 y_{p,\mathrm{lim}}=V_{DD}-y_{n,\mathrm{lim}}
 \approx\qty{0.3470}{\volt}.
\]
جای‌گذاری مستقیم در \eqref{eq:pass-p-body-implicit} این
مقدار را تأیید می‌کند.

\subsubsection{مسیر سری و کنترل گیت با سطح تضعیف‌شده}
\label{sec:pass-cascade}
شکل \secnum{۲.۳۱}\lr{(c)} کتاب و اسلاید ۱۳، چند \lr{nMOS}
سری با \lr{Gate}های مستقل متصل به \(V_{DD}\) را نشان
می‌دهند \cite{weste,dcslides}. در مدل آستانه ثابت، نخستین
افزاره حد \(V_{DD}-V_{tn}\) را تولید می‌کند.
افزاره بعدی با گیت کامل، از \eqref{eq:pass-n-min-rule} می‌دهد
\begin{equation}
 \min(V_{DD}-V_{tn},V_{DD}-V_{tn})=V_{DD}-V_{tn}.
 \label{eq:pass-series-same-gates}
\end{equation}
پس برای هر واحد بعدی با همان گیت کامل، حد انتقال همان است.
تعداد واحدها می‌تواند جریان، مقاومت مسیر و زمان شارژ را عوض
کند. تحلیل حاضر، حد ولتاژ بدون بار \lr{DC} را تعیین می‌کند.

در شکل \secnum{۲.۳۱}\lr{(d)}، خروجی تضعیف‌شده واحد اول به
\lr{Gate} واحد دوم وصل است. اگر ورودی مسیر دوم همچنان
\(V_{DD}\) باشد، گیت دوم فقط
\(V_{G,2}=V_{DD}-V_{tn}\) دارد. بنابراین،
\begin{equation}
 y_{2,\mathrm{lim}}=V_{G,2}-V_{tn}=V_{DD}-2V_{tn}.
 \label{eq:pass-gate-cascade}
\end{equation}
این نتیجه به وجود هدایت اولیه و مثبت‌بودن حد مورد نظر
وابسته است. افت بیشتر از کاهش ولتاژ کنترل گیت ناشی می‌شود
\cite[آغاز صفحه ۹۳]{weste}.
در مثال فرضی \(V_{DD}=\qty{1.2}{\volt}\) و
\(V_{tn}=\qty{0.3}{\volt}\)، مسیر سری با گیت‌های کامل حد
\(\qty{0.9}{\volt}\) دارد و ساختار کنترل گیت دوم با خروجی
اول، حد \(\qty{0.6}{\volt}\) می‌دهد.

افت سطح باید با مشخصات گیرنده مقایسه شود. اگر حد بالای
تولیدشده از \(V_{\mathrm{IH}}\) گیرنده کمتر باشد، سطح بالا
در همه شرایط تعریف‌شده معتبر نخواهد بود.
در تغذیه ثابت، نسبت \(V_t/V_{DD}\) سهم افت در دامنه منطقی
را تعیین می‌کند. کتاب، اهمیت آن را هنگام نزدیک‌شدن آستانه
به حدود یک‌سوم تغذیه در فناوری‌های مورد بحث خود توضیح
می‌دهد \cite[صفحه ۹۳]{weste}.

\subsubsection{عملکرد \lr{Transmission Gate}}
\label{sec:pass-transmission-gate}
\lr{Transmission Gate} یا \lr{TG} از یک \lr{nMOS} و یک
\lr{pMOS} موازی در مسیر سیگنال تشکیل می‌شود. دو گیت با
سیگنال‌های مکمل کنترل می‌شوند. در شکل
\ref{fig:pass-transmission-gate}، کنترل \(C=1\) یعنی
گیت \lr{nMOS} در \(V_{DD}\) و گیت \lr{pMOS} در زمین است.
برای \(C=0\)، دو ولتاژ کنترل جابه‌جا می‌شوند
\cite[صفحه ۹۳ و بخش \secnum{۱.۴.۶}]{weste}.
\begin{figure}[!htbp]
\centering
\begin{latin}
\PersianMathDigits
\begin{tikzpicture}[x=1cm,y=1cm]
\pic (tn) at (2,0.8) {study-npass};
\pic (tp) at (2,-0.8) {study-ppass};
\draw (tn-L)--(1,0.8)--(1,-0.8)--(tp-L);
\draw (tn-R)--(4,0.8)--(4,-0.8)--(tp-R);
\draw (0,0) node[ocirc] {} node[left] {\(V_{\mathrm{in}}\)}
 --(1,0) node[circ] {};
\draw (4,0) node[circ] {}--(5,0) node[ocirc] {}
 node[right] {\(V_{\mathrm{out}}\)};
\node[above] at (tn-G) {\(C\)};
\node[below] at (tp-G) {\(\overline C\)};
\node[right] at (3.2,1.15) {nMOS};
\node[right] at (3.2,-1.15) {pMOS};
\end{tikzpicture}
\end{latin}
\caption{ساختار آموزشی \lr{Transmission Gate} با کنترل مکمل.
دو افزاره در مسیر سیگنال موازی‌اند. بدنه \lr{nMOS} زمین و
بدنه \lr{pMOS} تغذیه است.}
\label{fig:pass-transmission-gate}
\end{figure}

در حالت روشن، \lr{nMOS} انتقال سطح پایین و \lr{pMOS}
انتقال سطح بالا را فراهم می‌کند. در مدل بدون بار \lr{DC}،
خروجی می‌تواند به ورودی برسد و افت آستانه تک‌افزاره‌ای در
دو ریل رفع می‌شود. در حالت خاموش، هر دو مسیر در مدل ایده‌آل
قطع می‌شوند و خروجی می‌تواند شناور باشد.

برای بررسی پوشش ولتاژ میانی با آستانه ثابت، فرض کنید دو
پایانه مسیر به ولتاژ مشترک \(v\) نزدیک‌اند. شرط روشن‌بودن
\lr{nMOS}، \(v<V_{DD}-V_{tn}\) و شرط روشن‌بودن
\lr{pMOS}، \(v>V_{TP}\) است. برای وجود پوشش هم‌پوشان
در سراسر بازه، شرط کافی این مدل چنین است
\begin{equation}
 V_{TP}<V_{DD}-V_{tn}
 \quad\Longleftrightarrow\quad V_{DD}>V_{tn}+V_{TP}.
 \label{eq:pass-tg-overlap}
\end{equation}
وقتی تغذیه پایین‌تر باشد، ممکن است بازه‌ای با هدایت ضعیف
برای هر دو افزاره ایجاد شود. با اثر بدنه، آستانه‌ها نیز
باید در بایاس همان گره بررسی شوند.

برای اختلاف ولتاژ کوچک دو پایانه، شیب جریان مدل مربعی
در نزدیکی \(V_{DS}=0\)، رسانایی روشن را می‌دهد. با
آستانه ثابت و تعریف \([x]_+=\max(x,0)\)، داریم
\begin{equation}
 g_{\mathrm{on}}(v)\approx
 \beta_n[V_{DD}-v-V_{tn}]_+
 +\beta_p[v-V_{TP}]_+,
 \qquad R_{\mathrm{on}}\approx\frac{1}{g_{\mathrm{on}}}.
 \label{eq:pass-tg-conductance}
\end{equation}
ضریب \(\beta\) با واحد آمپر بر ولت مربع در یک اضافه‌ولتاژ
ضرب می‌شود، پس رسانایی واحد زیمنس دارد. رسانایی دو مسیر
موازی جمع می‌شود. در مثال آموزشی
\(V_{DD}=\qty{1.2}{\volt}\)،
\(V_{tn}=V_{TP}=\qty{0.3}{\volt}\) و
\(\beta_n=\beta_p=\qty{200}{\micro\ampere\per\volt\squared}\)،
در \(v=\qty{0.6}{\volt}\) هر دو افزاره سهم
\(\qty{60}{\micro\siemens}\) دارند.
پس \(g_{\mathrm{on}}\approx\qty{120}{\micro\siemens}\) و
\(R_{\mathrm{on}}\approx\qty{8.33}{\kilo\ohm}\) است.
در نزدیکی هر ریل، سهم افزاره قوی‌تر حدود
\(\qty{180}{\micro\siemens}\) و مقاومت حدود
\(\qty{5.56}{\kilo\ohm}\) می‌شود.
این تقریب به اختلاف کوچک دو پایانه و مدل ساده جریان وابسته است.

\lr{TG} یک کلید کنترل‌شونده با ولتاژ است. در حالت روشن،
ولتاژ ورودی میانی را هم می‌تواند منتقل کند. بازیابی سطح
منطقی به ریل‌ها، به مشخصه یک گیت بازساز مانند
\lr{Inverter} وابسته است. مرجع کاربردهای \lr{TG} را در
\lr{Multiplexer}، ساختار منطقی، \lr{Latch} و کلید آنالوگ
ذکر می‌کند \cite[صفحه ۹۳]{weste}.

\subsubsection{حد \lr{DC}، بار خازنی و ورودی متغیر}
\label{sec:pass-dc-time}
در تحلیل حد انتقال، بار خروجی جریان ماندگار نمی‌گیرد.
اگر بار مقاومتی یا منبع جریان به خروجی وصل شود، باید
\lr{KCL} با جریان بار حل شود و مقدار خروجی می‌تواند تغییر
کند. در بار خازنی، جریان مسیر، ولتاژ را در زمان تغییر می‌دهد
و شرط \(I=0\) به حد نهایی مدل مربوط است.

اسلاید ۱۱، ورودی زمانی را معمولاً به شکل \lr{Step} یا
\lr{Ramp} معرفی می‌کند \cite{dcslides}.
\lr{Step} مدل تغییر ناگهانی از یک سطح به سطح دیگر است.
\lr{Ramp} تغییر خطی در یک بازه زمانی است. برای نمونه،
افزایش خطی از صفر به \(V_{DD}\) در مدت \(t_r>0\)،
رابطه زیر را در بازه مشخص‌شده دارد.
\[
 V_{\mathrm{in}}(t)=\frac{V_{DD}t}{t_r},\qquad 0\le t\le t_r.
\]
تحلیل زمانی به معادله تغییر بار و جریان وابسته به ولتاژ نیاز دارد.
نمونه‌ای از معادله گره خروجی چنین است.
\[
 C_L\frac{\mathrm{d}V_{\mathrm{out}}}{\mathrm{d}t}=I_{\mathrm{path}}-I_L.
\]
حد انتقال به‌تنهایی زمان
نشست را تعیین نمی‌کند.

\subsubsection{حل کامل تمرین \secnum{۲.۲۰}}
\label{sec:pass-ex20}
صورت اصلی تمرین، ولتاژ خروجی چهار مدار شکل
\secnum{۲.۳۵} را با صرف‌نظر از \lr{Body Effect} می‌خواهد
\cite[صفحه چاپی ۹۷]{weste}. شکل
\ref{fig:pass-ex20-circuits} بازترسیم همان اتصال‌هاست.
خط کوتاه افقیِ تغذیه به معنی \(V_{DD}\)، نماد زمین به معنی
صفر و حباب گیت مشخص‌کننده \lr{pMOS} است.
در هر مدار، گره میانی را \(X\) و خروجی را \(Y\) نامیده‌ایم.
\begin{figure}[!htbp]
\centering
\begin{latin}
\PersianMathDigits
\begin{tikzpicture}[x=1cm,y=1cm]
% Circuit (a).
\pic (aone) at (0,0) {study-npass};
\pic (atwo) at (2,0) {study-npass};
\draw (aone-L)--(-0.3,0) node[ground] {};
\draw (aone-R)--(atwo-L) node[midway,below] {\(X\)};
\draw (atwo-R)--(3.5,0) node[ocirc] {} node[right] {\(Y\)};
\draw (0.35,0.85)--(0.65,0.85) node[midway,above] {\(V_{DD}\)};
\draw (2.35,0.85)--(2.65,0.85) node[midway,above] {\(V_{DD}\)};
\node at (1.5,-1.65) {(a)};
% Circuit (b).
\begin{scope}[xshift=6.4cm]
 \pic (bone) at (0,0) {study-ppass};
 \pic (btwo) at (2,1.1) {study-ppass};
 \draw (bone-L)--(-0.3,0) node[ground] {};
 \draw (bone-G) node[ground] {};
 \draw (bone-R)--(2.5,0)--(btwo-G);
 \node[below] at (1.4,0) {\(X\)};
 \draw (btwo-L)--(1.7,1.1)--(1.7,0.7) node[ground] {};
 \draw (btwo-R)--(3.5,1.1) node[ocirc] {} node[right] {\(Y\)};
 \node at (1.5,-1.65) {(b)};
\end{scope}
% Circuit (c).
\begin{scope}[yshift=-4cm]
 \pic (cone) at (0,0) {study-ppass};
 \pic (ctwo) at (2,0) {study-npass};
 \draw (cone-L)--(-0.3,0) node[ground] {};
 \draw (cone-G) node[ground] {};
 \draw (cone-R)--(ctwo-L) node[midway,above] {\(X\)};
 \draw (ctwo-R)--(3.5,0) node[ocirc] {} node[right] {\(Y\)};
 \draw (2.35,0.85)--(2.65,0.85) node[midway,above] {\(V_{DD}\)};
 \node at (1.5,-1.65) {(c)};
\end{scope}
% Circuit (d).
\begin{scope}[xshift=6.4cm,yshift=-4cm]
 \pic (done) at (0,0) {study-ppass};
 \pic (dtwo) at (2,0) {study-npass};
 \draw (done-L)--(-0.3,0)--(-0.3,0.6);
 \draw (-0.45,0.6)--(-0.15,0.6) node[midway,above] {\(V_{DD}\)};
 \draw (done-G) node[ground] {};
 \draw (done-R)--(dtwo-L) node[midway,above] {\(X\)};
 \draw (dtwo-R)--(3.5,0) node[ocirc] {} node[right] {\(Y\)};
 \draw (2.35,0.85)--(2.65,0.85) node[midway,above] {\(V_{DD}\)};
 \node at (1.5,-1.65) {(d)};
\end{scope}
\end{tikzpicture}
\end{latin}
\caption{بازترسیم شکل \secnum{۲.۳۵} کتاب با نام‌گذاری گره‌ها.
در مدار \lr{(b)}، سیم \(X\) به گیت افزاره دوم وصل است.
نمادهای بدنه در شکل اصلی و این بازترسیم نمایش داده نشده‌اند.}
\label{fig:pass-ex20-circuits}
\end{figure}

مدل تمرین، آستانه ثابت و جریان بار و نشتی صفر دارد.
برای حالت‌های دارای افت، حدودی را محاسبه می‌کنیم که گره‌ها
از سمت قابل‌هدایت به آن‌ها نزدیک می‌شوند. در مدارهای
\lr{(b)} و \lr{(c)}، دشارژ از سطح بالاتر و در مدار
\lr{(d)}، شارژ خروجی از سطح پایین را در نظر می‌گیریم.
این فرض‌ها، تعبیر پاسخ‌های مدل \lr{Threshold Drop} را مشخص
می‌کنند. شرایط اولیه دلخواه گره‌های شناور، لزوماً همین
ولتاژها را ایجاد نمی‌کنند.

\paragraph*{مدار \lr{(a)} و عبور زمین}
ورودی مسیر صفر و گیت هر دو \lr{nMOS} در \(V_{DD}\) است.
با \(V_{DD}>V_{tn}\)، افزاره اول، \(X\) را تا صفر دشارژ
می‌کند. افزاره دوم نیز مسیر هدایت برای تخلیه \(Y\) به
همان صفر دارد. بنابراین،
\begin{equation}
 \boxed{V_X=0,\qquad V_Y=0}.
 \label{eq:pass-ex20-a}
\end{equation}
در حد نهایی، هر دو \(V_{DS}=0\) دارند و روشن‌بودن مسیر
با جریان صفر سازگار است. سری‌شدن دو واحد، افت آستانه‌ای
به سطح پایین اضافه نمی‌کند.

\paragraph*{مدار \lr{(b)} و کنترل گیت با سطح پایین تضعیف‌شده}
ترانزیستور اول \lr{pMOS} است. پایانه ورودی و گیت آن هر دو
زمین‌اند و گره \(X\) به گیت ترانزیستور دوم وصل است.
در دشارژ \(X\) از ولتاژ بالا، \(X\) پایانه با ولتاژ بیشتر
و در نتیجه سورس افزاره اول است. شرط هدایت
\(V_X-0>V_{TP}\) است. حد دشارژ گره کنترل برابر می‌شود با
\begin{equation}
 V_{X,\mathrm{lim}}=V_{TP}.
 \label{eq:pass-ex20-b-control}
\end{equation}
پایانه ورودی افزاره دوم نیز زمین است، اما گیت آن در
\(V_X\) قرار دارد. برای دشارژ خروجی از بالا، شرط هدایت
\(V_Y-V_X>V_{TP}\) است. با رسیدن گره کنترل به حد خود،
حد خروجی می‌شود
\begin{equation}
 \boxed{V_{Y,\mathrm{lim}}=V_X+V_{TP}=2V_{TP}}.
 \label{eq:pass-ex20-b}
\end{equation}
برای این تعبیر، \(V_{DD}>2V_{TP}\) را فرض می‌کنیم تا
خروجی بتواند از ولتاژ اولیه بالاتر به این حد دشارژ شود.
در حد نهایی، \(V_{SG}\) هر دو افزاره برابر \(V_{TP}\)
و جریان مدل صفر است. افت بیشتر به اتصال گیت دوم مربوط است.

برای روشن‌شدن نقش شرایط اولیه، اگر \(X\) از ابتدا صفر
باشد، افزاره اول در مدل بدون نشتی راهی برای شارژ آن تا
\(V_{TP}\) ندارد. گیت دوم در صفر می‌ماند و حد دشارژ
خروجیِ ابتدا بالا، \(V_{TP}\) می‌شود. پاسخ
\(2V_{TP}\)، حد متعارف مدل با دشارژ گره کنترل از بالا است.

\paragraph*{مدار \lr{(c)} و مسیر متشکل از \lr{pMOS} و \lr{nMOS}}
افزاره نخست، زمین را با گیت صفر به \(X\) منتقل می‌کند.
پس در دشارژ از بالا، \(V_{X,\mathrm{lim}}=V_{TP}\) است.
گیت \lr{nMOS} دوم در \(V_{DD}\) قرار دارد.
اگر \(V_{DD}-V_{TP}>V_{tn}\)، در نزدیکی حد نهایی،
این افزاره همچنان روشن است و اختلاف ولتاژ \(X\) و \(Y\)
را به صفر می‌رساند. بنابراین،
\begin{equation}
 \boxed{V_{Y,\mathrm{lim}}=V_{TP}}.
 \label{eq:pass-ex20-c}
\end{equation}
شرط تغذیه معادل \(V_{DD}>V_{TP}+V_{tn}\) است.
در حد نهایی، افزاره اول به آستانه رسیده و افزاره دوم با
اضافه‌ولتاژ مثبت، اختلاف درین و سورس صفر دارد.
در این مدار، \lr{nMOS} همان سطح پایین تضعیف‌شده را منتقل
می‌کند و افت جداگانه‌ای به آن نمی‌افزاید.

\paragraph*{مدار \lr{(d)} و انتقال تغذیه}
ورودی \lr{pMOS} اول \(V_{DD}\) و گیت آن صفر است.
با \(V_{DD}>V_{TP}\)، این افزاره \(X\) را تا
\(V_{DD}\) شارژ می‌کند. ورودی و گیت \lr{nMOS} دوم در
این حد هر دو \(V_{DD}\) هستند. در شارژ \(Y\) از پایین،
اضافه‌ولتاژ آن \(V_{DD}-V_Y-V_{tn}\) است. پس
\begin{equation}
 \boxed{V_X=V_{DD},\qquad V_{Y,\mathrm{lim}}=V_{DD}-V_{tn}}.
 \label{eq:pass-ex20-d}
\end{equation}
در این حد، \lr{pMOS} روشن و اختلاف دو پایانه کانال آن صفر
است. \lr{nMOS} به مرز آستانه رسیده است.

جدول \ref{tab:pass-ex20-results} پاسخ‌ها و شرط‌های آن‌ها را
جمع می‌کند. برای مقایسه عددی آموزشی، اگر
\(V_{DD}=\qty{1}{\volt}\) و
\(V_{tn}=V_{TP}=\qty{0.3}{\volt}\) باشند، حدهای خروجی
\lr{(a)} تا \lr{(d)} به‌ترتیب صفر،
\(\qty{0.6}{\volt}\)، \(\qty{0.3}{\volt}\) و
\(\qty{0.7}{\volt}\) خواهند بود. داده عددی در صورت اصلی
مسئله تعیین نشده است.
\begin{table}[!htbp]
\centering
\small
\caption{نتایج تمرین \secnum{۲.۲۰} در مدل آستانه ثابت و با مسیر
شارژ یا دشارژ توضیح‌داده‌شده.}
\label{tab:pass-ex20-results}
\begin{tabularx}{\linewidth}{cc>{\raggedleft\arraybackslash}X}
\toprule
مدار & خروجی یا حد انتقال & شرط تغذیه برای تحلیل\\
\midrule
\lr{(a)} & صفر & \(V_{DD}>V_{tn}\)\\
\lr{(b)} & \(2V_{TP}\) & \(V_{DD}>2V_{TP}\)\\
\lr{(c)} & \(V_{TP}\) & \(V_{DD}>V_{TP}+V_{tn}\)\\
\lr{(d)} & \(V_{DD}-V_{tn}\) & \(V_{DD}>\max(V_{TP},V_{tn})\)\\
\bottomrule
\end{tabularx}
\end{table}
این جواب‌ها، حدود مدل بدون نشتی را بیان می‌کنند. با اثر
بدنه یا بار \lr{DC}، آستانه‌ها و معادلات جریان باید متناسب
با بایاس واقعی هر افزاره به‌روز شوند.

\clearpage
\BookletReview{notes/booklet-03/final-review}
